% !TEX root = ../main.tex
\documentclass[12pt]{report}
\usepackage[a4paper,left=0.885in,right=0.885in,top=1.33in,bottom=1.46in,includefoot]{geometry}
\usepackage{fontspec}
% Times New Roman on Windows; TeX Gyre Termes (Times-compatible) on Overleaf/Linux.
\IfFontExistsTF{Times New Roman}{%
  \setmainfont{Times New Roman}%
}{%
  \setmainfont{TeX Gyre Termes}%
}
% 한글 글꼴(xetexko). 본문은 명조, 산세리프는 고딕, 고정폭은 고딕코딩.
% Linux·Overleaf에는 나눔 글꼴, Windows에는 바탕·맑은 고딕을 씁니다.
\usepackage{kotex}
\IfFontExistsTF{NanumMyeongjo}{%
  \setmainhangulfont{NanumMyeongjo}[AutoFakeSlant=0.15]%
}{%
  \IfFontExistsTF{Batang}{%
    \setmainhangulfont{Batang}[AutoFakeSlant=0.15]%
  }{%
    \setmainhangulfont{Malgun Gothic}[AutoFakeSlant=0.15]%
  }%
}
\IfFontExistsTF{NanumGothic}{%
  \setsanshangulfont{NanumGothic}%
}{%
  \setsanshangulfont{Malgun Gothic}%
}
\IfFontExistsTF{NanumGothicCoding}{%
  \setmonohangulfont{NanumGothicCoding}[AutoFakeSlant=0.15]%
}{%
  \setmonohangulfont{Malgun Gothic}%
}
\usepackage{setspace}
\usepackage{hyperref}
\usepackage{xurl}
\usepackage{seqsplit}
\usepackage{microtype}
\usepackage{fvextra}
\usepackage{enumitem}
\usepackage{tikz}
\usetikzlibrary{arrows.meta,positioning,shapes.geometric,fit,calc,decorations.pathreplacing,patterns}
\usepackage{longtable}
\usepackage{booktabs}
\usepackage{siunitx}
\sisetup{round-mode=places,round-precision=3,detect-weight=true,detect-family=true}
\usepackage{array}
\usepackage[normalem]{ulem}
\usepackage{caption}
\usepackage{amsmath}
\usepackage{amssymb}
\usepackage{amsthm}
\usepackage{graphicx}
\usepackage{float}
\usepackage{algorithm}
\usepackage{algpseudocode}
\usepackage{multirow}
\usepackage{tabularx}
\usepackage{adjustbox}
\usepackage{subcaption}
\usepackage[bottom,hang,flushmargin]{footmisc}
% Keep tables within the text block (no right-margin overflow).
\newcommand{\fitwidth}[1]{%
  \begin{adjustbox}{max width=\linewidth,center}%
  \small\setlength{\tabcolsep}{3.5pt}%
  #1%
  \end{adjustbox}%
}
\hypersetup{
  hidelinks,
  pdftitle={ERC-20 허락 화살표를 페이지랭크에 넣어 온체인 지갑 평판 점수를 더 좋게 만들기 (한국어 쉬운 번역)},
  pdfauthor={권태홍 (Taehong Kwon)},
  pdfsubject={남아프리카대학교(UNISA) 컴퓨터과학 박사 학위 논문의 한국어 번역},
  pdfkeywords={DeFi, 엔도스랭크, EndorseRank, 페이지랭크, 온체인 평판, Arbitrum}
}
% 한국어 이름표: 그림, 표, 목록, 알고리즘
\AtBeginDocument{%
  \renewcommand{\figurename}{그림}%
  \renewcommand{\tablename}{표}%
  \renewcommand{\listfigurename}{그림 목록}%
  \renewcommand{\listtablename}{표 목록}%
  \renewcommand{\contentsname}{차례}%
  \renewcommand{\chaptername}{장}%
  \renewcommand{\appendixname}{부록}%
}
\floatname{algorithm}{알고리즘}
\algrenewcommand\algorithmicrequire{\textbf{입력:}}
\algrenewcommand\algorithmicensure{\textbf{출력:}}
\algrenewcommand\algorithmicreturn{\textbf{돌려주기:}}
\algrenewtext{For}[1]{\textbf{반복:} #1}
\algrenewtext{EndFor}{\textbf{반복 끝}}
\def\UrlFont{\ttfamily}
\emergencystretch=2em
\DeclareRobustCommand{\texttt}[1]{{\ttfamily\seqsplit{#1}}}
\RecustomVerbatimEnvironment{verbatim}{Verbatim}{%
  breaklines=true,breakanywhere=true,fontsize=\small,formatcom=\ttfamily}
\setstretch{1.58}
\setlength{\parindent}{0pt}
\setlength{\parskip}{0.03\baselineskip}
\setlist[itemize]{leftmargin=2.2em,labelsep=0.55em,itemsep=0.08\baselineskip,topsep=0.1\baselineskip,parsep=0pt}
\setlist[enumerate]{leftmargin=2.4em,labelsep=0.55em,itemsep=0.08\baselineskip,topsep=0.1\baselineskip,parsep=0pt}
\usepackage{titlesec}
\usepackage{etoolbox}
\pretocmd{\chapter}{\phantomsection}{}{}
% Match 1-Proposal article \section* / \subsection* sizing and spacing.
\titleformat{\chapter}[block]
  {\normalfont\Large\bfseries}{\thechapter}{0.5em}{}
\titlespacing*{\chapter}{0pt}{3.5ex plus 1ex minus .2ex}{2.3ex plus .2ex}
\titleformat{name=\chapter,numberless}[block]
  {\normalfont\Large\bfseries}{}{0pt}{}
\titlespacing*{name=\chapter,numberless}{0pt}{3.5ex plus 1ex minus .2ex}{2.3ex plus .2ex}
\titleformat{\section}{\normalfont\large\bfseries}{\thesection}{0.5em}{}
\titlespacing*{\section}{0pt}{3.25ex plus 1ex minus .2ex}{1.5ex plus .2ex}
\titleformat{\subsection}{\normalfont\normalsize\bfseries}{\thesubsection}{0.5em}{}
\titlespacing*{\subsection}{0pt}{2.5ex plus 1ex minus .2ex}{1.2ex plus .2ex}
\setcounter{secnumdepth}{3}
\setcounter{tocdepth}{2}
\newcommand{\dissertationSubheading}[2][]{%
  \phantomsection
  \section{#2}%
  \if\relax\detokenize{#1}\relax\else\label{#1}\fi
}
\newcommand{\dissertationSubsubheading}[2][]{%
  \phantomsection
  \subsection{#2}%
  \if\relax\detokenize{#1}\relax\else\label{#1}\fi
}
\newcommand{\TOCentry}[3]{%
  \noindent\hyperref[#3]{\makebox[2.4em][l]{#1}#2}\nobreak\dotfill\ \hyperref[#3]{\pageref*{#3}}\par\vspace{0.15em}%
}
\newcommand{\dissertationResultHeading}[1]{%
  \par\addvspace{0.45\baselineskip}%
  \noindent\textbf{#1}\par\nobreak\addvspace{0.06\baselineskip}%
}
\newcommand{\figpath}[1]{Figures/generated/#1}

% Citations: superscript [n] in text, hyperlinked to Section 7 References.
% No page-bottom citation footnotes (媛곸＜); full entries live only in References.
\ExplSyntaxOn
\cs_new_protected:Npn \cend_one:n #1
  {
    \textsuperscript{\hyperref[ref:#1]{[#1]}}
  }
\cs_new_protected:Npn \cend:n #1
  {
    \clist_set:Nn \l_tmpa_clist { #1 }
    \clist_sort:Nn \l_tmpa_clist
      {
        \int_compare:nNnTF { ##1 } > { ##2 }
          { \sort_return_swapped: }
          { \sort_return_same: }
      }
    \clist_clear:N \l_tmpb_clist
    \clist_map_inline:Nn \l_tmpa_clist
      {
        \clist_if_in:NnF \l_tmpb_clist { ##1 }
          { \clist_put_right:Nn \l_tmpb_clist { ##1 } }
      }
    \cs_set_nopar:Npn \cend_sep: {}
    \clist_map_inline:Nn \l_tmpb_clist
      {
        \cend_sep:
        \cend_one:n { ##1 }
        \cs_set_nopar:Npn \cend_sep: { \textsuperscript{,} }
      }
  }
\cs_new_eq:NN \cend \cend:n
\ExplSyntaxOff

